% Esquema de un viaducto de n vanos y longitud total L: la luz media es L/n si los vanos son iguales.
\documentclass[tikz,border=4pt]{standalone}
\usepackage{resonancia}
\begin{document}
\begin{tikzpicture}
  % tablero y pilas (6 vanos iguales)
  \draw[grueso, color=rGris] (0,1.0) -- (12,1.0);
  \foreach \x in {0,2,4,6,8,10,12} {
    \draw[grueso, color=rGris] (\x,1.0) -- (\x,0.2);
    \fill[apoyo] (\x-0.25,0) rectangle (\x+0.25,0.2);
    \draw[trazo, color=rGris] (\x-0.25,0.2) -- (\x+0.25,0.2);
  }
  % cota de un vano
  \draw[<->, trazo, color=rAzul] (4,1.45) -- (6,1.45);
  \node[etiqueta, color=rAzul, anchor=south] at (5,1.62) {luz de un vano};
  \draw[punteada] (4,1.05) -- (4,1.55);
  \draw[punteada] (6,1.05) -- (6,1.55);
  % cota total
  \draw[<->, trazo, color=rNaranja] (0,-0.5) -- (12,-0.5);
  \node[etiqueta, color=rNaranja, anchor=north] at (6,-0.55) {longitud total $L$ (suma de los $n$ vanos)};
  \draw[punteada] (0,-0.05) -- (0,-0.5);
  \draw[punteada] (12,-0.05) -- (12,-0.5);
  % fórmula
  \node[etiqueta, anchor=west, align=left] at (12.6,0.7) {luz media\\[2pt]$\bar{l}=\dfrac{L}{n}$\\[2pt]\scriptsize vanos iguales};
\end{tikzpicture}
\end{document}
